%% kapa: exam=1 round=36 year=2025 period=1 form=A subject=real_estate subject_ko=부동산학원론 seq=6 number=86
%% source: Q-Net 감정평가사 기출문제 (cst00309) artlSeq=5231525 file=raw/1008/5231525/2025년 제36회 감정평가사 1차시험 1교시 문제지 원본.pdf
%% conversion: kordoc 4.24.51 -> tools/gen1.py (auto); review=proofread-sonnet; flags=none
\begin{question}{86}
부동산의 특성에 관한 설명으로 옳지 않은 것은?
\begin{choices}
  \item 부동산은 물리적 구성요소들의 결합체일 뿐만 아니라 여러가지 경제적ㆍ사회적 특성의 결합체이다.
  \item 토지의 개별성으로 인해 일물일가의 법칙이 적용되지 않고, 부동산상품 간에 완전한 대체는 불가능하다.
  \item 토지거래허가구역의 지정이나 해제 등으로 인해 주택가격이 하락하거나 상승하는 것은 경제적 위치의 변화에 따른 것이다.
  \item 토지는 용도전환 및 합병ㆍ분할이 가능하며, 두 개 이상의 용도가 병존할 수도 있다.
  \item 토지의 부동성으로 인해 부동산시장을 국지화시키며, 이로 인해 부동산의 가치는 주변 환경의 영향을 많이 받는다.
\end{choices}
\end{question}
\answer{3}
